\section{A. Dinamika Pelaksanaan Undang-Undang Dasar Negara Republik Indonesia Tahun 1945 (UUD NRI Tahun 1945)}

UUD NRI Tahun 1945 merupakan sumber hukum tertinggi di Indonesia. Pembahasan mengenai konstitusi selalu berkaitan dengan Pancasila sebagai sumber dari segala sumber hukum (\textit{staatsfundamentalnorm}). Ditetapkan oleh PPKI pada tanggal 18 Agustus 1945.

\subsection{Wawasan Kewarganegaraan}
\begin{itemize}
    \item UUD NRI Tahun 1945 yang berlaku sekarang ditetapkan berdasarkan Dekrit Presiden 5 Juli 1959 beserta perubahan-perubahannya (amandemen).
    \item UUD NRI Tahun 1945 merupakan pengganti UUDS 1950.
\end{itemize}

\subsection{1. Pentingnya Konstitusi}
\begin{itemize}
    \item Istilah konstitusi berasal dari bahasa Prancis (\textit{constituer}) dan Inggris (\textit{constitution}) yang berarti membentuk.
    \item Makna kata konstitusi selalu terikat dalam konteks pembentukan suatu organ dan organisasi negara.
    \item Dibedakan menjadi konstitusi tertulis (\textit{documentary constitution}) dan konstitusi tidak tertulis (\textit{non-documentary constitution} / konvensi).
\end{itemize}

Konstitusi adalah segala ketentuan dan aturan mengenai ketatanegaraan atau Undang-Undang Dasar suatu negara. Negara dan konstitusi merupakan dua hal yang tidak dapat dipisahkan. Konstitusi merupakan peraturan atau hukum dasar tertinggi yang dijadikan pegangan dalam penyelenggaraan pemerintahan negara. Terdapat dua jenis konstitusi, yaitu tertulis dan tidak tertulis.

\subsubsection{Wawasan Kewarganegaraan: Tiga Unsur Konstitusi}
Menurut Savonir Lohman (dalam Solly Lubis, 1982), tubuh konstitusi memuat tiga unsur:
\begin{enumerate}
    \item Perwujudan perjanjian masyarakat (kontrak sosial) untuk membina negara dan pemerintahan yang mengatur mereka.
    \item Piagam penjaminan hak-hak asasi manusia serta penentuan batas hak dan kewajiban warga negara maupun alat pemerintahan.
    \item \textit{Forma regimenis} (kerangka bangunan atau gambaran struktur pemerintahan negara).
\end{enumerate}

\subsubsection{Muatan Pokok Konstitusi}
Menurut J.G. Steenbeek, terdapat tiga muatan pokok konstitusi:
\begin{itemize}
    \item Pengaturan tentang perlindungan hak asasi manusia dan warga negara.
    \item Pengaturan tentang susunan ketatanegaraan suatu negara yang mendasar.
    \item Pembatasan dan pembagian tugas-tugas ketatanegaraan yang mendasar.
\end{itemize}

Konstitusi berfungsi membatasi wewenang dan kekuasaan agar tidak dilakukan menurut keinginan penguasa secara sewenang-wenang (konstitusionalisme). UUD merupakan bagian dari konstitusi tertulis (\textit{documentary constitution}) yang dilengkapi dengan peraturan/kebiasaan tidak tertulis (konvensi) dan yurisprudensi.

\subsubsection{Contoh Kebiasaan Ketatanegaraan (Konvensi) di Indonesia}
\begin{itemize}
    \item Pidato Presiden setiap tanggal 16 Agustus di depan sidang paripurna DPR.
    \item Praktik pengambilan keputusan di MPR berdasarkan musyawarah mufakat.
    \item Presiden menyampaikan penjelasan RUU APBN di minggu pertama bulan Januari di hadapan DPR.
    \item Keberadaan Menteri Negara nondepartemen pada masa pemerintahan Orde Baru.
\end{itemize}

\subsubsection{Wawasan Kewarganegaraan: Pidato Presiden Tanggal 16 Agustus}
Sejak tahun 2010, tradisi ketatanegaraan ini diformalkan melalui UU No. 27 Tahun 2009 (UU MD3). Presiden melakukan dua pidato dalam sehari pada tanggal 16 Agustus: Pidato Kenegaraan menyambut HUT Kemerdekaan RI (dalam Sidang Bersama DPR dan DPD) dan Pidato Penyampaian RUU APBN (di hadapan DPR).


\subsection{2. Sejarah Penyusunan Undang-Undang Dasar Negara Republik Indonesia Tahun 1945 (UUD NRI Tahun 1945)}
UUD NRI Tahun 1945 dirancang sejak tanggal 29 Mei 1945 sampai 16 Juli 1945 oleh BPUPK (\textit{Dokuritsu Junbi Coosakai}). Hasil rumusan diawali dari Piagam Jakarta (22 Juni 1945) dan dibahas pada masa persidangan kedua BPUPK (10--17 Juli 1945). Menurut Bernhard Dahm (1987), terjadi silang pendapat antara kubu nasionalis dan agamis mengenai bentuk negara (negara kebangsaan atau negara Islam).

\subsubsection{Wawasan Kewarganegaraan: Sidang BPUPK Kedua}
Sidang diikuti oleh 74 orang dipimpin oleh Ketua KRT Radjiman Wedyodiningrat serta Wakil Ketua Ichibangase Yosio dan R. Panji Soeroso.

Tiga masalah pokok rancangan UUD yang disampaikan Sukarno pada Sidang BPUPK kedua:
\begin{enumerate}
    \item Pernyataan tentang Indonesia Merdeka.
    \item Pembukaan Undang-Undang Dasar.
    \item Batang Tubuh Undang-Undang Dasar.
\end{enumerate}

\subsubsection{Pembentukan Panitia Kerja BPUPK}
Pada tanggal 10 Juli 1945, disepakati bentuk negara kesatuan dan dibentuk tiga kelompok panitia:
\begin{itemize}
    \item \textbf{Panitia Perancang Hukum Dasar}: Diketuai oleh Ir. Sukarno.
    \item \textbf{Panitia Perancang Keuangan dan Ekonomi}: Diketuai oleh Drs. Mohammad Hatta.
    \item \textbf{Panitia Perancang Pembelaan Tanah Air}: Diketuai oleh Abikoesno Tjokrosoejoso.
\end{itemize}

\subsubsection{Wawasan Kewarganegaraan: Tabel 2.1 Rancangan Undang-Undang Dasar}
\begin{table}[h!]
\centering
\begin{tabular}{|l|l|p{8cm}|}
\hline
\textbf{No.} & \textbf{Tanggal Sidang} & \textbf{Hasil Rancangan} \\ \hline
1. & 10 Juli 1945 & Keputusan bentuk negara dan pembentukan tiga kelompok panitia. \\ \hline
2. & 11--12 Juli 1945 & Pokok-Pokok Pikiran Pembukaan UUD dan Rancangan UUD (Hasil Rancangan Panitia Kecil). \\ \hline
3. & 13 Juli 1945 & Rancangan Pertama UUD dibahas pada forum Panitia Perancang Hukum Dasar. \\ \hline
4. & 14 Juli 1945 & Rancangan Kedua UUD dibahas pada forum Sidang BPUPK. \\ \hline
5. & 15--16 Juli 1945 & Rancangan Ketiga UUD (terakhir) disahkan dalam Sidang BPUPK. \\ \hline
\end{tabular}
\end{table}

\subsubsection{Panitia Kecil dan Panitia Penghalus Bahasa}
Panitia Perancang Hukum Dasar membentuk Panitia Kecil (7 orang) pada 11 Juli 1945 diketuai Prof. Dr. Mr. Supomo. Selain itu, dibentuk Panitia Penghalus Bahasa yang terdiri atas Husein Djajadiningrat, Agoes Salim, dan Prof. Dr. Mr. Supomo.

\subsubsection{Pembentukan PPKI dan Pengesahan UUD 1945}
\begin{itemize}
    \item BPUPK dibubarkan pada 7 Agustus 1945, lalu pada 9 Agustus 1945 dibentuk Panitia Persiapan Kemerdekaan Indonesia (PPKI) dengan 21 anggota.
    \item Proklamasi Kemerdekaan dilaksanakan pada 17 Agustus 1945.
    \item Pada sidang 18 Agustus 1945, PPKI mengesahkan UUD NRI Tahun 1945, memilih Ir. Sukarno sebagai Presiden dan Drs. Mohammad Hatta sebagai Wakil Presiden, serta membentuk Komite Nasional.
    \item Tanggal 18 Agustus ditetapkan sebagai Hari Konstitusi berdasarkan Keputusan Presiden Nomor 18 Tahun 2008.
\end{itemize}

\subsubsection{Tokoh: Kasman Singodimedjo}
Kasman Singodimedjo (lahir 25 Februari 1904) merupakan Komandan PETA Jakarta, pengaman Proklamasi 17 Agustus 1945, anggota tambahan PPKI, Ketua KNIP pertama, Jaksa Agung, dan pahlawan nasional (Kepres No. 123/TK/Tahun 2018). Beliau berperan sebagai pemersatu bangsa dalam menjembatani dialog antara golongan Islam dan nasionalis mengenai penyempurnaan tujuh kata Piagam Jakarta pada 18 Agustus 1945.


\subsection{3. Pembukaan Undang-Undang Dasar Negara Republik Indonesia Tahun 1945 (UUD NRI Tahun 1945)}
Pembukaan UUD NRI Tahun 1945 terdiri atas empat alinea dan berkedudukan di atas pasal-pasal UUD sebagai *staatsfundamentalnorm* (norma hukum dasar tertinggi) yang memiliki hubungan kausal organis.

\subsubsection{a. Isi dan Tujuan Pembukaan Undang-Undang Dasar Negara Republik Indonesia Tahun 1945 (UUD NRI Tahun 1945)}
Menurut Kaelan (2020), makna per alinea Pembukaan UUD NRI 1945 adalah:
\begin{enumerate}
    \item \textbf{Alinea Pertama}: Mengandung pengakuan hak kodrat bahwa kemerdekaan adalah hak segala bangsa. Penjajahan bertentangan dengan perikemanusiaan dan perikeadilan.
    \item \textbf{Alinea Kedua}: Konsekuensi logis dari pernyataan kemerdekaan dan bukti objektif atas perjuangan bangsa menuju cita-cita negara yang merdeka, bersatu, berdaulat, adil, dan makmur.
    \item \textbf{Alinea Ketiga}: Pengakuan nilai religius dan moral bahwa kemerdekaan merupakan rahmat Tuhan Yang Maha Kuasa dan hak kodrat.
    \item \textbf{Alinea Keempat}: Memuat tujuan khusus dan tujuan umum negara, ketentuan diadakannya UUD Negara, bentuk negara Republik yang berkedaulatan rakyat, serta dasar filsafat negara (Pancasila).
\end{enumerate}

\subsubsection{Wawasan Kewarganegaraan: Negara yang Merdeka, Bersatu, dan Berdaulat}
\begin{itemize}
    \item \textbf{Merdeka}: Bebas dari kekuasaan dan campur tangan bangsa lain.
    \item \textbf{Bersatu}: Mengandung arti kebulatan dan kesatuan bangsa Indonesia untuk mewujudkan negara persatuan.
    \item \textbf{Berdaulat}: Berdiri di atas kemampuan, kekuatan, dan kekuasaan sendiri serta memiliki kedudukan sederajat dengan bangsa lain.
\end{itemize}

\subsubsection{Makna yang terkandung dalam Pembukaan UUD NRI Tahun 1945, yaitu sebagai berikut}
\begin{enumerate}
    \item Sumber dari motivasi dan aspirasi tekad dan perjuangan bangsa Indonesia.
    \item Sumber dari cita-cita hukum dan moral dalam lingkup nasional dan pergaulan antarbangsa.
    \item Mempunyai nilai-nilai universal (dijunjung tinggi oleh bangsa-bangsa di dunia) dan lestari (menampung dinamika masyarakat).
\end{enumerate}

Terdapat empat poin penegasan alinea dalam Pembukaan UUD NRI Tahun 1945:
\begin{enumerate}
    \item Penegasan pernyataan kemerdekaan adalah hak segala bangsa dan penjajahan harus dihapuskan.
    \item Penetapan cita-cita bangsa Indonesia yang hendak dicapai dengan kemerdekaan.
    \item Penegasan Proklamasi Kemerdekaan 17 Agustus 1945 sebagai rahmat Tuhan Yang Maha Kuasa.
    \item Penegasan ketentuan negara yang meliputi tujuan dan bentuk negara, ketentuan UUD, dan rumusan filsafat dasar Negara Indonesia.
\end{enumerate}

\subsubsection{b. Pembukaan UUD NRI Tahun 1945 sebagai Pokok Kaidah Negara yang Fundamental (\textit{Staatsfundamentalnorm})}

\subsubsection{Dua Unsur Utama \textit{Staatsfundamentalnorm}}
\begin{enumerate}
    \item \textbf{Dari Unsur Terjadinya}: Dibentuk oleh pembentuk negara (PPKI) sebagai wakil bangsa Indonesia dan penjelang dasar negara.
    \item \textbf{Dari Unsur Isinya}: Memuat dasar tujuan negara, ketentuan diadakannya UUD, bentuk negara, dan dasar filsafat negara (Pancasila).
\end{enumerate}

\subsubsection{Alasan Pembukaan UUD NRI 1945 Tidak Dapat Diubah}
\begin{itemize}
    \item Ditetapkan oleh PPKI yang mewakili seluruh bangsa Indonesia. Setelah PPKI dibubarkan, tidak ada lagi lembaga/alat negara yang berhak mengubahnya.
    \item Mengubah atau menghapus Pembukaan sama artinya dengan \textbf{membubarkan Negara Kesatuan Republik Indonesia (NKRI)}.
    \item Secara formal-yuridis, kedudukannya bersifat tetap, kuat, dan terlekat erat pada keberadaan negara.
\end{itemize}


\subsubsection{c. Hubungan Pembukaan UUD NRI Tahun 1945 dengan Proklamasi Kemerdekaan 17 Agustus 1945}
Pembukaan UUD NRI Tahun 1945 dan Proklamasi Kemerdekaan merupakan \textbf{satu kesatuan yang tidak dapat dipisahkan}.

\begin{itemize}
    \item \textbf{Alinea III} memuat pernyataan kemerdekaan yang menyatu dengan spirit teks Proklamasi.
    \item Pengesahan Pembukaan UUD 1945 pada 18 Agustus 1945 merupakan \textbf{penjelasan, penegasan, dan pertanggungjawaban} pelaksanaan Proklamasi Kemerdekaan.
    \item Proklamasi berfungsi sebagai pintu gerbang/jembatan emas kemerdekaan, sedangkan Pembukaan memberikan penegasan arah, bentuk, dan tujuan kehidupan bernegara.
\end{itemize}


\subsubsection{d. Hubungan Pembukaan UUD NRI Tahun 1945 dengan Pasal-Pasalnya}
\begin{itemize}
    \item Pembukaan memuat norma-norma dasar (\textit{staatsfundamentalnorm}), sedangkan pasal-pasal merupakan penjabaran teknis dari norma dasar tersebut ke dalam hukum tertulis.
    \item \textbf{Empat Pokok Pikiran dalam Pembukaan UUD NRI Tahun 1945}:
    \begin{enumerate}
        \item \textbf{Pokok Pikiran Persatuan}: Negara melindungi segenap bangsa Indonesia dan seluruh tumpah darah Indonesia.
        \item \textbf{Pokok Pikiran Keadilan Sosial}: Negara hendak mewujudkan keadilan sosial bagi seluruh rakyat Indonesia.
        \item \textbf{Pokok Pikiran Kedaulatan Rakyat}: Negara berkedaulatan rakyat berdasarkan kerakyatan dan permusyawaratan/perwakilan.
        \item \textbf{Pokok Pikiran Ketuhanan \& Kemanusiaan}: Negara berdasarkan Ketuhanan Yang Maha Esa menurut dasar kemanusiaan yang adil dan beradab.
    \end{enumerate}
\end{itemize}